% Paso 5: Análisis de Género Literario
% Proyecto: Análisis Exegético Profundo
% Ministerio "La Gloria es del Señor"

\documentclass[11pt,a4paper]{article}
\usepackage[utf8]{inputenc}
\usepackage[spanish]{babel}
\usepackage[margin=2.5cm]{geometry}
\usepackage{graphicx}
\usepackage{fancyhdr}
\usepackage{xcolor}
\usepackage{tcolorbox}
\usepackage{enumitem}
\usepackage{tabularx}
\usepackage{multirow}
\usepackage{amsfonts}
\usepackage{fontspec}
\usepackage{titlesec}

% Configuración de colores
\definecolor{forestgreen}{RGB}{27,67,50}
\definecolor{sagegreen}{RGB}{45,90,39}
\definecolor{olivegreen}{RGB}{64,145,108}
\definecolor{mintgreen}{RGB}{82,183,136}
\definecolor{lightmint}{RGB}{116,198,157}
\definecolor{softcream}{RGB}{216,243,220}
\definecolor{papergray}{RGB}{118,128,150}

% Configuración de encabezados
\pagestyle{fancy}
\fancyhf{}
\fancyhead[L]{\textcolor{forestgreen}{\textbf{Análisis Exegético Profundo}}}
\fancyhead[R]{\textcolor{papergray}{Paso 5: Género Literario}}
\fancyfoot[C]{\textcolor{papergray}{\thepage}}
\fancyfoot[R]{\textcolor{papergray}{\scriptsize Ministerio "La Gloria es del Señor"}}

% Configuración de títulos
\titleformat{\section}
{\color{forestgreen}\Large\bfseries}
{\thesection}{1em}{}

\titleformat{\subsection}
{\color{sagegreen}\large\bfseries}
{\thesubsection}{1em}{}

% Configuración de cajas
\tcbuselibrary{skins}
\newtcolorbox{infobox}[1]{
    colback=softcream,
    colframe=olivegreen,
    fonttitle=\bfseries,
    title=#1,
    left=10pt,
    right=10pt,
    top=10pt,
    bottom=10pt
}

\newtcolorbox{checklistbox}{
    colback=white,
    colframe=mintgreen,
    left=15pt,
    right=15pt,
    top=10pt,
    bottom=10pt,
    boxrule=2pt
}

\newtcolorbox{genrebox}{
    colback=lightmint!20,
    colframe=sagegreen,
    left=10pt,
    right=10pt,
    top=8pt,
    bottom=8pt,
    boxrule=1.5pt
}

\newtcolorbox{literarybox}{
    colback=olivegreen!10,
    colframe=olivegreen,
    left=10pt,
    right=10pt,
    top=8pt,
    bottom=8pt,
    boxrule=1.5pt
}

\begin{document}

% Título principal
\begin{center}
    {\Huge\color{forestgreen}\textbf{PASO 5}}\\[0.5cm]
    {\Large\color{sagegreen}\textbf{GÉNERO LITERARIO}}\\[0.3cm]
    {\large\color{papergray}Análisis de Forma y Estructura Literaria}\\[1cm]
    
    \begin{tcolorbox}[colback=olivegreen!10, colframe=olivegreen, width=0.8\textwidth]
        \centering
        \textcolor{forestgreen}{\textbf{Proyecto: Análisis Exegético Profundo}}\\
        \textcolor{papergray}{Metodología de 8 Pasos para Estudio Bíblico Riguroso}
    \end{tcolorbox}
\end{center}

\vspace{1cm}

% Información del estudiante
\begin{infobox}{Información del Proyecto}
\begin{tabularx}{\textwidth}{|X|X|}
\hline
\textbf{Nombre del Estudiante:} & \rule{6cm}{0.4pt} \\[0.8cm]
\hline
\textbf{Texto Bíblico:} & \rule{6cm}{0.4pt} \\[0.8cm]
\hline
\textbf{Libro Bíblico:} & \rule{3cm}{0.4pt} \quad \textbf{Testamento:} \rule{2cm}{0.4pt} \\[0.8cm]
\hline
\textbf{Fecha:} & \rule{6cm}{0.4pt} \\[0.8cm]
\hline
\end{tabularx}
\end{infobox}

\section{Objetivo del Paso 5: Género Literario}

El análisis del género literario es crucial para la interpretación correcta. Cada género tiene sus propias convenciones, propósitos y métodos de comunicación. Identificar y entender el género nos ayuda a aplicar los principios hermenéuticos apropiados.

\subsection{Descripción del Proceso}

Identificaremos el género literario específico y sus características distintivas. Aplicaremos principios hermenéuticos apropiados para narrativa, poesía, epístola, profecía, etc., y analizaremos la estructura literaria del pasaje.

\begin{infobox}{Géneros Bíblicos Principales}
\textbf{Narrativa, Ley, Poesía, Sabiduría, Profecía, Evangelio, Epístola, Apocalíptica}\\
Cada uno requiere enfoques interpretativos específicos y únicos.
\end{infobox}

\section{Lista de Verificación}

\begin{checklistbox}
\textbf{Marca cada elemento completado:}

\begin{itemize}[leftmargin=1cm]
    \item[$\square$] Clasificar el género literario principal
    \item[$\square$] Identificar características específicas del género
    \item[$\square$] Aplicar principios hermenéuticos apropiados
    \item[$\square$] Considerar elementos retóricos y estilísticos
    \item[$\square$] Analizar la estructura literaria del pasaje
    \item[$\square$] Identificar subgéneros o formas literarias
    \item[$\square$] Examinar el propósito comunicativo del género
    \item[$\square$] Evaluar paralelos en literatura contemporánea
\end{itemize}
\end{checklistbox}

\newpage

\section{Identificación del Género}

\subsection{Clasificación Primaria}

\begin{genrebox}
\textbf{Selecciona el género literario principal de tu texto:}
\end{genrebox}

\begin{itemize}[leftmargin=1cm]
    \item[$\square$] \textbf{Narrativa Histórica} - Relatos de eventos del pasado
    \item[$\square$] \textbf{Ley/Legal} - Códigos legales, mandamientos, ordenanzas
    \item[$\square$] \textbf{Poesía} - Salmos, cantos, literatura poética
    \item[$\square$] \textbf{Literatura Sapiencial} - Proverbios, reflexiones sobre la vida
    \item[$\square$] \textbf{Profecía} - Oráculos, visiones, mensajes proféticos
    \item[$\square$] \textbf{Evangelio} - Narrativa sobre la vida de Jesús
    \item[$\square$] \textbf{Epístola} - Cartas apostólicas, correspondencia
    \item[$\square$] \textbf{Apocalíptica} - Literatura de revelación/visión
    \item[$\square$] \textbf{Parábola} - Relatos ilustrativos con significado espiritual
    \item[$\square$] \textbf{Otro:} \rule{8cm}{0.4pt}
\end{itemize}

\subsection{Evidencia para la Clasificación}

\textbf{¿Qué características del texto apoyan esta clasificación de género?}
\vspace{3cm}

\textbf{¿Hay elementos de otros géneros presentes?}
\vspace{2cm}

\section{Características del Género}

\subsection{Elementos Distintivos}

\begin{tabularx}{\textwidth}{|X|X|}
\hline
\textbf{Característica del Género} & \textbf{Evidencia en el Texto} \\
\hline
Propósito típico del género & \\[1.5cm]
\hline
Estructura común & \\[1.5cm]
\hline
Estilo y lenguaje & \\[1.5cm]
\hline
Audiencia típica & \\[1.5cm]
\hline
Método de comunicación & \\[1.5cm]
\hline
\end{tabularx}

\newpage

\section{Análisis por Género Específico}

\subsection{Para Narrativa}

\begin{literarybox}
\textbf{Completa solo si tu texto es narrativo:}
\end{literarybox}

\textbf{Elementos de la narrativa presentes:}
\begin{itemize}[leftmargin=1cm]
    \item[$\square$] Personajes claramente definidos
    \item[$\square$] Trama con conflicto y resolución
    \item[$\square$] Ambiente/escenario establecido
    \item[$\square$] Punto de vista narrativo consistente
    \item[$\square$] Tema(s) central(es)
\end{itemize}

\textbf{Estructura narrativa:}
\begin{enumerate}[leftmargin=1cm]
    \item Exposición: \rule{10cm}{0.4pt}
    \item Conflicto/Tensión: \rule{10cm}{0.4pt}
    \item Clímax: \rule{10cm}{0.4pt}
    \item Resolución: \rule{10cm}{0.4pt}
\end{enumerate}

\textbf{¿Es narrativa descriptiva o prescriptiva?}
\vspace{1.5cm}

\subsection{Para Poesía}

\begin{literarybox}
\textbf{Completa solo si tu texto es poético:}
\end{literarybox}

\textbf{Tipos de paralelismo identificados:}
\begin{itemize}[leftmargin=1cm]
    \item[$\square$] Sinónimo (repite la idea)
    \item[$\square$] Antitético (contrasta ideas)
    \item[$\square$] Sintético (expande la idea)
    \item[$\square$] Quiástico (estructura cruzada)
    \item[$\square$] Escalado (construcción progresiva)
\end{itemize}

\textbf{Recursos poéticos presentes:}
\begin{itemize}[leftmargin=1cm]
    \item[$\square$] Metáforas
    \item[$\square$] Símiles
    \item[$\square$] Personificación
    \item[$\square$] Hipérbole
    \item[$\square$] Aliteración
    \item[$\square$] Inclusio (marco literario)
\end{itemize}

\subsection{Para Epístola}

\begin{literarybox}
\textbf{Completa solo si tu texto es epistolar:}
\end{literarybox}

\textbf{Elementos epistolares presentes:}
\begin{itemize}[leftmargin=1cm]
    \item[$\square$] Saludo/Identificación del remitente
    \item[$\square$] Acción de gracias
    \item[$\square$] Cuerpo principal
    \item[$\square$] Exhortaciones prácticas
    \item[$\square$] Saludos finales/Bendición
\end{itemize}

\textbf{Tipo de sección epistolar:}
\begin{itemize}[leftmargin=1cm]
    \item[$\square$] Doctrinal/Teológica
    \item[$\square$] Práctica/Ética
    \item[$\square$] Correctiva/Disciplinaria
    \item[$\square$] Personal/Pastoral
\end{itemize}

\newpage

\subsection{Para Profecía}

\begin{literarybox}
\textbf{Completa solo si tu texto es profético:}
\end{literarybox}

\textbf{Tipos de material profético:}
\begin{itemize}[leftmargin=1cm]
    \item[$\square$] Oráculo de juicio
    \item[$\square$] Oráculo de salvación
    \item[$\square$] Lamento
    \item[$\square$] Himno
    \item[$\square$] Visión simbólica
    \item[$\square$] Acción simbólica
\end{itemize}

\textbf{Estructura del oráculo:}
\begin{enumerate}[leftmargin=1cm]
    \item Fórmula del mensajero: \rule{8cm}{0.4pt}
    \item Acusación: \rule{8cm}{0.4pt}
    \item Anuncio de juicio/salvación: \rule{8cm}{0.4pt}
\end{enumerate}

\textbf{Tiempo profético:}
\begin{itemize}[leftmargin=1cm]
    \item[$\square$] Cercano (cumplimiento inmediato)
    \item[$\square$] Lejano (cumplimiento escatológico)
    \item[$\square$] Doble cumplimiento
\end{itemize}

\section{Estructura Literaria}

\subsection{Organización del Material}

\textbf{¿Cómo está estructurado tu pasaje?}

\begin{tabularx}{\textwidth}{|c|X|X|}
\hline
\textbf{Sección} & \textbf{Versículos} & \textbf{Contenido/Función} \\
\hline
A & & \\[1cm]
\hline
B & & \\[1cm]
\hline
C & & \\[1cm]
\hline
D & & \\[1cm]
\hline
E & & \\[1cm]
\hline
\end{tabularx}

\subsection{Patrones Estructurales}

\textbf{¿Hay patrones estructurales evidentes?}
\begin{itemize}[leftmargin=1cm]
    \item[$\square$] Quiasmo (A-B-C-B'-A')
    \item[$\square$] Inclusio (marco literario)
    \item[$\square$] Paralelismo estructural
    \item[$\square$] Progresión temática
    \item[$\square$] Estructura concéntrica
    \item[$\square$] Otro: \rule{6cm}{0.4pt}
\end{itemize}

\textbf{Diagrama de la estructura (si es compleja):}
\vspace{4cm}

\newpage

\section{Elementos Retóricos}

\subsection{Estrategias Retóricas}

\textbf{¿Qué técnicas retóricas usa el autor para comunicar su mensaje?}

\begin{tabularx}{\textwidth}{|X|X|}
\hline
\textbf{Técnica Retórica} & \textbf{Ejemplo en el Texto} \\
\hline
Repetición & \\[1cm]
\hline
Contraste & \\[1cm]
\hline
Pregunta retórica & \\[1cm]
\hline
Ironía & \\[1cm]
\hline
Clímax & \\[1cm]
\hline
Énfasis por posición & \\[1cm]
\hline
\end{tabularx}

\subsection{Propósito Retórico}

\textbf{¿Cuál es el efecto retórico pretendido en la audiencia original?}
\vspace{3cm}

\textbf{¿Cómo contribuye el género a lograr este propósito?}
\vspace{2cm}

\section{Principios Hermenéuticos}

\subsection{Principios Específicos del Género}

\textbf{¿Qué principios hermenéuticos son especialmente relevantes para este género?}
\vspace{3cm}

\textbf{¿Qué errores interpretativos comunes debes evitar con este género?}
\vspace{3cm}

\subsection{Aplicación Apropiada}

\textbf{¿Cómo afecta el género la forma en que debes aplicar este texto?}
\vspace{3cm}

\section{Síntesis del Análisis}

\subsection{Hallazgos Clave}

\textbf{¿Qué aspectos del análisis de género más iluminan tu texto?}
\vspace{3cm}

\textbf{¿Cómo afecta esta comprensión del género tu interpretación general?}
\vspace{3cm}

\textbf{¿Hay algún elemento de género que cambió tu comprensión inicial del texto?}
\vspace{3cm}

\begin{infobox}{Principio Hermenéutico}
El género no determina completamente el significado, pero sí establece las "reglas del juego" para la interpretación. Respeta las convenciones del género mientras buscas el mensaje del autor.
\end{infobox}

\section{Siguiente Paso}

Con el análisis de género completado, procederás al \textbf{Paso 6: Contexto Canónico}, donde examinaremos el texto dentro del contexto de toda la Escritura.

\vspace{1cm}

\begin{center}
\begin{tcolorbox}[colback=mintgreen!20, colframe=mintgreen, width=0.7\textwidth]
\centering
\textcolor{forestgreen}{\textbf{¡Análisis Literario Completo!}}\\
\textcolor{papergray}{Has identificado las características de género que guían la interpretación.}
\end{tcolorbox>
\end{center>

\end{document}